% !PRINTED 649
\weberpage{649}
%<<CONT

Haben zwei holomorphe Funktionen
\[
\begin{aligned}
\alpha &= \alpha_0 x^h + \alpha_1 x^{h-1} + \cdots + \alpha_h, \\
\beta &= \beta_0 x^k + \beta_1 x^{k-1} + \cdots + \beta_k
\end{aligned}
\]
ein Produkt
\[
\gamma = \gamma_0 x^{h+k} + \gamma_1 x^{h+k-1} + \cdots + \gamma_{h+k},
\]
dessen Koeffizienten ganze Funktionale sind, so können die $\gamma_0$,
$\gamma_1, \ldots, \gamma_{h+k}$ nur dann alle durch ein Primfunktional $\pi$ teilbar
sein, wenn entweder alle $\alpha_0, \alpha_1, \ldots, \alpha_h$, oder alle $\beta_0, \beta_1, \ldots, \beta_k$
durch $\pi$ teilbar sind, was ganz so bewiesen wird, wie in Bd. I, § 2.

Wenn nun (2) ein ganzes Funktional ist, ohne daß $\varphi_0, \varphi_1, \ldots, \varphi_m$
ganz sind, so kann man ein ganzes Funktional $\mu$ und darin einen
Primfaktor $\pi$ so bestimmen, daß
\weq{3}{\chi = \mu\varphi = \alpha_0 x^m + \alpha_1 x^{m-1} + \cdots + \alpha_m,}
durch $\pi$ teilbar ist, $\alpha_0, \alpha_1, \ldots, \alpha_n$ aber nicht alle durch $\pi$ teilbar
% [sic] printed $\alpha_n$, the sum in (3) runs to $\alpha_m$
sind.

Ist nun $\varphi$ ganz, so genügt es einer Gleichung:
\weq{4}{E\varphi^m = a_1 E_1 \varphi^{m-1} + a_2 E_2 \varphi^{m-2} + \cdots + a_m E_m}
und daraus durch Multiplikation mit $\mu^m$:
\weq{5}{E\chi^m = \mu(a_1 E_1 \chi^{m-1} + a_2 E_2 \mu \chi^{m-2} + a_m E_m),}
% [sic] printed $\mu(a_1 E_1 \chi^{m-1} + a_2 E_2 \mu \chi^{m-2} + a_m E_m)$, multiplying (4) by $\mu^m$ gives $\mu(a_1 E_1 \chi^{m-1} + a_2 E_2 \mu \chi^{m-2} + \cdots + a_m E_m \mu^{m-1})$; the $+ \cdots$ and the factor $\mu^{m-1}$ of the last term are missing
worin die $a_1, a_2, \ldots$ ganze Funktionen in $Z$, die $E, E_1, \ldots, E_m$
Einheiten in $\bar{Z}$ sind. Die Einheiten $E, E_1, \ldots, E_m$ können zwar
$x$ noch enthalten, können aber holomorph angenommen werden.
Ordnen wir rechter Hand und linker Hand von (5) nach $x$, so
sind die Koeffizienten von $E$ nicht alle durch $\pi$ teilbar. Ebenso
sind nach dem erwähnten Satz die Koeffizienten von $\chi^m$ nicht
alle durch $\pi$ teilbar, während auf der rechten Seite alle Koeffizienten
 den Faktor $\mu$, also auch den Faktor $\pi$ haben. Das ist
unmöglich und damit unser Satz bewiesen (Bd. II, § 159).

8. \gesperrt{Ein holomorphes ganzes Funktional $\varphi$ ist der
größte gemeinschaftliche Teiler aller seiner Koeffizienten.}

Sind $\alpha, \beta, \ldots$ die Koeffizienten von $\varphi$, so ist $\varphi$ jedenfalls
durch den größten gemeinschaftlichen Teiler von $\alpha, \beta, \ldots$ teilbar.
Zerlegen wir eine Linearfunktion $z - c$ in ihre funktionalen Primfaktoren,
so können wir in diesen Primfunktionalen $\pi$ die Variablen
beliebig bezeichnen, und da jedes Primfunktional in einem $z - c$
aufgehen muß, so können wir darin die Variablen von den Variablen
 $x, y, \ldots$ in $\varphi$ verschieden annehmen. Demnach können
%CONT>>
